% !TeX root = ../main-pinn.tex
% !TeX program = xelatex
\chapter{物理信息神经网络与风场重建}
\label{chap:pinn}
\begin{quote}\small
\textit{本章摘要。}激光雷达只能测到少量位置、特定方向上的风速，我们却希望知道一片区域中的风怎样变化。物理信息神经网络（Physics-Informed Neural Network，PINN）把测量数据和流体方程一起用于训练，为这一问题提供了一种求解方法。本章先通过手算例子说明训练过程，介绍 PINN 的通用框架与正、逆问题，再讨论雷达观测、风场方程和损失函数，结合具体论文理解风场重建，最后介绍训练与模型的优化方向。
\end{quote}

想象风机前方有一片随时间变化的风场。激光雷达沿几条波束测量风速，得到的是分散的读数。而风机控制和风场分析希望知道区域内各处的速度和方向。根据这些有限读数补出连续风场，就是本章要讨论的任务。

如果只按测量值做插值，结果主要取决于测点的位置和插值规则。流体力学还提供了另一类信息：相邻位置、相邻时刻的速度应当遵循质量守恒和动量变化规律。PINN 用神经网络表示风场，同时用数据和方程来调整这个网络\cite{karniadakis2021physics, raissi2019pinn}。

本章以二维风场为主，先把方法的计算过程讲清楚。这里的二维风场指某个平面内的两个速度分量，记为$u(x,y,t)$和$v(x,y,t)$。学会这个例子后，再把空间坐标和速度分量扩展到三维，会更容易理解。

\section{先用一个小例子理解 PINN}
\label{sec:pinn-formulation}

\subsection{网络表示的是一个可以查询的函数}
先把要求解的未知函数记为$q$。如果它只有一个自变量$x$，就可以让网络接收$x$，输出这个位置的函数值：

\begin{equation}
 q_\theta(x)=F_\theta(x).
 \label{eq:pinn-generic-network}
\end{equation}
下标$\theta$表示网络中需要学习的权重和偏置。训练就是调整这些参数，使网络表示的函数接近我们要求的解。训练完成后，给定一个$x$，就能查询对应的$q_\theta(x)$。

如果未知函数还随时间变化，输入就增加时间$t$。如果要同时求解几个物理量，输出就增加相应分量。先从一个只有$x$的例子开始，便于看清方程怎样参与训练。

\subsection{把方程变成可以计算的误差}
先暂时离开风场，考虑一个只有一个自变量的方程：
\begin{equation}
 q'(x)+q(x)=0,\qquad q(0)=1,\qquad 0\leq x\leq1.
 \label{eq:pinn-hand-ode}
\end{equation}
这个方程的解是$q(x)=e^{-x}$。这里假设我们还不知道解，想通过训练找到它。

为方便手算，先用一条直线$q_{a,b}(x)=ax+b$近似这个函数，$a,b$就是待学习参数。它对坐标的导数是$q'_{a,b}=a$。把直线代入方程左侧，得到
\[
 r(x)=q'_{a,b}(x)+q_{a,b}(x)=a+ax+b.
\]
这个$r(x)$叫作\textbf{方程残差}，表示当前函数与方程要求的差距。我们希望它在区间内尽量接近零。

先选$x=0,1$两个位置检查方程，再检查$q(0)=1$这个条件。将误差平方后相加，得到损失函数
\begin{equation}
 \mathcal L(a,b)=\frac12\big[(a+b)^2+(2a+b)^2\big]+(b-1)^2.
 \label{eq:pinn-hand-loss}
\end{equation}
前面两项检查方程，最后一项检查$x=0$处的函数值。像$x=0,1$这样用来计算方程残差的位置，称为\textbf{配点}。选择配点只需要坐标，这些位置不必有真实函数值标签。

从$a=0,b=1$开始，当前函数是$q(x)=1$。它在$x=0$处满足给定条件，但两个位置的方程残差都是$1$，总损失为$1$。此时
\[
 \frac{\partial\mathcal L}{\partial a}=3,\qquad
 \frac{\partial\mathcal L}{\partial b}=2.
\]
取学习率$0.1$，沿减小损失的方向更新一次：
\[
 a\leftarrow0-0.1\times3=-0.3,\qquad
 b\leftarrow1-0.1\times2=0.8.
\]
新函数为$q(x)=-0.3x+0.8$。两个位置的残差变为$0.5$和$0.2$，总损失下降到
\[
 \frac12(0.5^2+0.2^2)+(0.8-1)^2=0.185.
\]
这就完成了一次训练：先计算函数和导数，再算残差与损失，最后更新参数。实际 PINN 用多层神经网络代替这条直线，重复同样的过程，逐渐得到更合适的函数。

\subsection{两次求导分别在做什么}
上面的计算中出现了两种导数。对坐标$x$求导，得到$q'(x)$，用来检查方程。对参数$a,b$求导，得到损失的下降方向，用来更新模型。风场 PINN 也是这样：先计算速度对位置、时间的导数，再计算损失对网络参数的导数。

\textbf{自动微分}就是由程序按照链式法则完成这些求导。只要网络的计算过程由可求导的运算组成，就可以从输出追踪到输入，计算$\partial_xu$、$\partial_tu$等量。可以先把它理解为训练框架提供的一套求导工具，后续实现时再学习具体接口。

\subsection{同一套流程可以用于哪些问题}
\label{sec:pinn-generality}
手算例子中的关键步骤是：用一个可学习函数表示未知量，把它代入方程形成残差，再结合已知条件计算损失、更新参数。PINN 的通用性就体现在这套流程可以用于不同的方程和物理系统\cite{raissi2019pinn,wang2023expert}。

换一个任务时，先回答三个问题：\textbf{输入哪些坐标，输出哪些未知量，用什么方程检查它们？}表\ref{tab:pinn-general-problems}列出几个例子。稳态问题可以只输入空间坐标，随时间演化的问题则加入$t$。输出既可以是一个温度，也可以是多个速度分量和压力。

\begin{table}[htbp]
\centering\small
\caption{不同任务中的 PINN 输入、输出与方程。这里以一维传热、波动和二维流动为例。}
\label{tab:pinn-general-problems}
\begin{tabular}{m{0.18\linewidth}m{0.13\linewidth}m{0.17\linewidth}m{0.37\linewidth}}
\toprule
任务 & 输入 & 输出 & 用于训练的方程\\
\midrule
本章手算例子 & $x$ & $q$ & $q'+q=0$\\
一维热传导 & $(x,t)$ & 温度 $T$ & $\partial_tT-\kappa\partial_{xx}T=0$\\
一维波动 & $(x,t)$ & 位移 $s$ & $\partial_{tt}s-c^2\partial_{xx}s=0$\\
二维风场 & $(x,y,t)$ & $u,v,\pi$ & 连续性与动量方程\\
\bottomrule
\end{tabular}
\end{table}

例如，把未知量换成温度$T(x,t)$，建立网络$T_\theta=F_\theta(x,t)$，再按热传导方程计算
\begin{equation}
 r_T(x,t)=\partial_tT_\theta-\kappa\partial_{xx}T_\theta.
 \label{eq:pinn-generic-heat-residual}
\end{equation}
$\kappa$是热扩散系数。这里的$r_T$与手算例子的$r=q'+q$承担相同作用，都是检查当前函数与方程要求的差距。热传导中还可以给出初始温度、两端的温度或热流条件，把这些条件也写成误差项。

更一般地，可以用$\mathcal N[q;\mu]=0$概括要满足的方程。其中，$\mathcal N$表示方程中的求导和组合运算，$\mu$表示方程参数，例如热扩散系数$\kappa$或波速$c$。确定任务后，就按相应方程构造残差，结合已知条件和可用观测形成损失。网络训练得到的是在所选表示、配点和损失权重下的近似解。小的离散残差不能单独保证全域解误差小，也不能保证方程参数可由当前观测唯一确定。

\subsection{正问题、逆问题与本章的风场重建}
\paragraph{正问题：已知模型与条件，求物理场。}
例如已知材料的热扩散系数、初始温度和两端条件，希望求出之后各处的温度。训练时使用给定方程和这些条件，在内部配点上检查残差，最后查询网络得到温度分布。这类任务称为正问题。

\paragraph{逆问题：利用观测，推断未知场或参数。}
例如已经测到少量位置的温度，希望估计材料的热扩散系数。可以将$\kappa$与网络权重一起作为待学习参数，用测量误差和方程残差共同训练。观测来自哪里、测到了什么，就通过相应的观测关系写入损失\cite{raissi2019pinn}。

本章关心的是另一类逆问题：从少量雷达读数恢复未测位置的速度场。网络以时空坐标为输入，一种常见的输出形式是
\begin{equation}
 (u_\theta,v_\theta,\pi_\theta)=F_\theta(x,y,t).
 \label{eq:pinn-coordinate-network}
\end{equation}
$u,v$是两个速度分量，$\pi=p/\rho$是压力除以空气密度得到的量。压力与速度变化通过动量方程联系，因此也可以作为网络需要学习的未知场。测量数据通过训练时的误差参与学习，训练后的网络则可以在整个区域内查询。

从这里开始，我们把通用框架中的未知函数换成风场，依次弄清雷达测到了什么、流体方程怎样写、两类信息怎样一起训练网络。

\section{激光雷达给了我们什么数据}
\label{sec:pinn-lidar}

\subsection{径向速度是风速在波束方向上的投影}
二维速度记为$\mathbf u=(u,v)^\top$。设一条波束与$x$轴的夹角为$\alpha$，方向向量是$\mathbf e=(\cos\alpha,\sin\alpha)^\top$，理想点测量可写成
\begin{equation}
 y_{\mathrm{LOS}}=\mathbf e^\top\mathbf u
 =u\cos\alpha+v\sin\alpha.
 \label{eq:pinn-observation}
\end{equation}
LOS 是 line of sight 的缩写，表示视线方向。测量值也称径向速度。

例如某处的速度为$(u,v)=(3,4)\,\mathrm{m/s}$。沿$x$方向的波束测到$3\,\mathrm{m/s}$，沿$y$方向的波束测到$4\,\mathrm{m/s}$，而风速大小为$\sqrt{3^2+4^2}=5\,\mathrm{m/s}$。建立数据损失时，要先把网络预测的速度投影到波束方向，再与仪器读数比较。

若同一位置、同一时刻有两个不平行方向的测量，就能联立两个投影关系求解$u,v$。实际扫描中的不同波束通常经过不同位置，此时可以借助流体方程把分散的观测联系起来。后面的论文案例展示的就是这种做法。

\subsection{一条训练样本包含哪些信息}
一条二维雷达观测可以整理为
\[
 (x_m,y_m,t_m,\alpha_m,y_m^{\mathrm{LOS}}).
\]
它依次记录测点位置、测量时间、波束方向和径向速度。训练时，将$(x_m,y_m,t_m)$送进网络得到$u_\theta,v_\theta$，再用$\alpha_m$计算预测的径向速度。

可以把投影计算统一写成$\widehat y_m=\mathcal H_m[\mathbf u_\theta]$。符号$\mathcal H_m$叫作\textbf{观测算子}，可以理解为“把预测风场转换成第$m$条仪器读数的计算规则”。对于式\eqref{eq:pinn-observation}，它就是一次投影。

实际雷达的一个距离门会覆盖一定长度，测量过程还可能包含时间平均。需要描述这些效应时，可以在若干位置和时刻查询网络，再加权求和：
\begin{equation}
 \widehat y_m=\sum_{a=1}^{A_m}\omega_{ma}\,
 \mathbf e_{ma}^{\top}\mathbf u_\theta(\mathbf x_{ma},t_{ma}),
 \qquad \sum_{a=1}^{A_m}\omega_{ma}=1.
\end{equation}
$A_m$是求和使用的点数，$\omega_{ma}$是权重。第一次学习时，可以先掌握点投影，再根据仪器的距离门设置扩展这一步\cite{sathe2013lidar}。

数据整理时，先统一坐标轴、角度定义、速度正负号和时间单位。例如仪器方位角可能从北方向起算，而模型中的角度从$x$轴起算，转换后才能使用同一条投影公式。

\section{用流体方程描述风速怎样变化}
\label{sec:pinn-physics}

\subsection{连续性方程与局部流量平衡}
先考虑恒密度、不可压缩的二维流动，其连续性方程为
\begin{equation}
 \frac{\partial u}{\partial x}+\frac{\partial v}{\partial y}=0.
 \label{eq:pinn-continuity}
\end{equation}
可以在风场中想象一个很小的方框。空气从方框的一边进入，从另一边流出。在密度近似不变的情况下，各边流量需要平衡。式\eqref{eq:pinn-continuity}表达的就是这种局部平衡。

其中$\partial_xu$描述$u$沿$x$方向的变化率，$\partial_yv$描述$v$沿$y$方向的变化率。将网络预测代入方程左侧，得到连续性残差
\[
 r_c=\partial_xu_\theta+\partial_yv_\theta.
\]
训练时让$r_c$接近零，就在推动预测风场满足这一平衡关系。

\subsection{动量方程中的速度变化与受力}
连续性方程回答的是“空气的流入和流出怎样平衡”。接下来还要回答：\textbf{空气为什么会加速、减速或改变方向？}这就需要动量方程。

先把一小团空气想成一个正在运动的物体。牛顿第二定律告诉我们，物体受到合力就会产生加速度。对空气也是如此：周围的压力和黏性作用会改变它的速度。把牛顿第二定律用于流体，并除以这团空气的质量，就得到“加速度等于单位质量所受的合力”这样的关系。

在本章恒密度、不可压缩的二维模型中，暂不加入重力等体积力，Navier--Stokes 动量方程可以展开为
\begin{align}
 \partial_tu+u\partial_xu+v\partial_yu
 &=-\partial_x\pi+\nu(\partial_{xx}u+\partial_{yy}u),\\
 \partial_tv+u\partial_xv+v\partial_yv
 &=-\partial_y\pi+\nu(\partial_{xx}v+\partial_{yy}v).
 \label{eq:pinn-momentum}
\end{align}
这里$\nu$是运动黏度，单位为$\mathrm{m^2/s}$。$\pi=p/\rho$是单位质量压力，量纲为$\mathrm{m^2/s^2}$，其中$p$是压力，$\rho$是恒定空气密度。两行方程分别描述$x$、$y$方向的运动。该模型假设二维、不可压缩、恒黏度流动，并暂不计重力等体积力。先读第一行，可以把它概括为
\[
 \underbrace{\partial_tu+u\partial_xu+v\partial_yu}_{\text{跟着空气团看到的加速度}}
 =\underbrace{-\partial_x\pi}_{\text{压力作用}}
 +\underbrace{\nu(\partial_{xx}u+\partial_{yy}u)}_{\text{黏性作用}}.
\]
等号两边的单位都是$\mathrm{m/s^2}$。下面逐步说明这三个部分。

\paragraph{左边：为什么加速度要写成三项？}
$u(x,y,t)$描述每个位置、每个时刻的$x$方向速度。其中$\partial_tu$表示固定位置的速度随时间变化有多快，$\partial_xu$和$\partial_yu$则表示同一时刻不同位置之间的速度差异。前者像把风速仪固定在一点，连续记录读数。后两者像在不同位置放置风速仪，比较同一时刻的读数。

如果跟着一团空气移动，就要同时考虑时间和位置的变化。在很短的$\Delta t$内，它沿$x$方向移动约$u\Delta t$，沿$y$方向移动约$v\Delta t$。因此，它的$x$方向速度变化可以近似写成
\[
 \Delta u\approx\partial_tu\,\Delta t
 +\partial_xu\,(u\Delta t)+\partial_yu\,(v\Delta t).
\]
两边除以$\Delta t$，再让时间间隔趋近于零，就得到第一行左侧的三个加速度项。它们合起来称为$u$的\textbf{物质导数}，常记为$Du/Dt$。其中$u\partial_xu+v\partial_yu$称为\textbf{对流项}，记录空气移动到新位置时，空间速度差异带来的变化。

图\ref{fig:pinn-material-acceleration}展示了一种容易混淆的情形：各个固定位置的风速都不随时间变化，但越往右风速越大。固定的风速仪读数不变，沿流动方向移动的空气团却会逐渐变快。例如，在某一点$u=4\,\mathrm{m/s}$、$\partial_xu=0.05\,\mathrm{s^{-1}}$，仅$u\partial_xu$这一项就贡献了$0.2\,\mathrm{m/s^2}$的加速度。这解释了为什么只用$\partial_tu$还不足以描述空气团的加速度。

\begin{figure}[htbp]
\centering
\begin{tikzpicture}[font=\small]
 \node at (4,1.3) {固定位置的读数不变，沿 $x$ 方向的速度逐渐增大};
 \draw[->] (-0.5,-0.15)--(9.4,-0.15) node[right] {$x$};
 \foreach \xx/\speed/\len in {0/4/1.0,3.5/5/1.4,7/6/1.8} {
  \fill (\xx,0.45) circle (1.6pt);
  \draw[idi arrow,blue!65!black] (\xx,0.45)--++(\len,0);
  \node[above] at (\xx+0.5,0.6) {$u=\speed\,\mathrm{m/s}$};
 }
 \node[below] at (0,-0.15) {$x_A$};
 \node[below] at (3.5,-0.15) {$x_B$};
 \node[below] at (7,-0.15) {$x_C$};
 \draw[idi feedback] (0,-1.0)--node[below,text=black]{跟着空气团向右移动，速度逐渐增大}(7,-1.0);
\end{tikzpicture}
\caption{固定位置的变化与跟随空气团的变化。实线箭头表示同一时刻各处的$x$方向速度，黑点标出测量位置。虚线箭头表示空气团的移动方向，数字为教学示例。}
\label{fig:pinn-material-acceleration}
\end{figure}

\paragraph{右边第一项：压力差怎样推动空气？}
可以想象这团空气的左、右两侧都受到周围空气的挤压。如果左侧压力较高、右侧较低，向右的推动就更强，空气会获得向右的加速度。此时压力沿$x$方向下降，$\partial_xp<0$，所以$-\partial_x\pi>0$，正好表示向右的作用。除以密度，是把单位体积的压力作用换成单位质量对应的加速度。

\paragraph{右边第二项：黏性为什么与速度的二阶导数有关？}
相邻空气层的速度不同时，会通过黏性作用交换动量。较快的空气层会带动较慢的层，自己也会受到减速作用，整体效果是使速度差异逐渐变得平缓。$\nu$表示这种动量扩散的强弱。

二阶导数$\partial_{xx}u+\partial_{yy}u$描述速度分布在两个方向上的弯曲程度。例如，只看沿$y$方向的速度剖面，在一个局部速度峰顶附近，$\partial_{yy}u<0$，黏性项倾向于降低这个峰值。在局部低谷附近，$\partial_{yy}u>0$，它则倾向于抬高低谷。这样就能把“平滑速度差异”与公式中的二阶导数联系起来。

第二行按同样的方式理解：跟着空气团观察$y$方向速度$v$的变化，右边由$y$方向的压力差和黏性作用提供加速度。\textbf{连续性方程检查流量平衡，两个动量方程检查速度变化与受力是否匹配。}这三条方程共同构成本章用于训练风场的物理关系。

和手算例子一样，把所有项移到左边，就能得到两个动量残差：
\begin{align}
 r_u&=\partial_tu_\theta+u_\theta\partial_xu_\theta
 +v_\theta\partial_yu_\theta+\partial_x\pi_\theta
 -\nu(\partial_{xx}u_\theta+\partial_{yy}u_\theta),\\
 r_v&=\partial_tv_\theta+u_\theta\partial_xv_\theta
 +v_\theta\partial_yv_\theta+\partial_y\pi_\theta
 -\nu(\partial_{xx}v_\theta+\partial_{yy}v_\theta).
 \label{eq:pinn-dimensional-residual}
\end{align}
三种残差都可以在区域内部的配点上计算。观测点负责告诉网络“这里测到了多少”，配点则负责检查“这里的变化是否符合方程”。

\subsection{初值、边界与尺度}
\textbf{初始条件}描述开始时刻的状态，例如$t=0$时的速度分布。\textbf{边界条件}描述区域边缘的状态，例如入口风速或地表条件。在实际重建中，先整理能获得的入口、初始场等信息，再把它们写成相应误差项。论文案例采用的具体损失会在第\ref{sec:pinn-validation}节说明。

风场中的坐标可能是几百米，时间是几十秒，速度是几米每秒。训练前通常先选择长度尺度$L_0$和速度尺度$U_0$，把各变量换成相对大小：
\begin{align}
 x^*=\frac{x}{L_0},\quad y^*=\frac{y}{L_0},\quad
 t^*=\frac{tU_0}{L_0},\qquad
 u^*=\frac{u}{U_0},\quad v^*=\frac{v}{U_0},\quad
 \pi^*=\frac{\pi}{U_0^2}.
 \label{eq:pinn-scales}
\end{align}
这一步称为\textbf{无量纲化}。例如$U_0=10\,\mathrm{m/s}$时，$u=3\,\mathrm{m/s}$变成$u^*=0.3$。训练时预测和观测使用同一套尺度，输出结果时再换回米、秒和米每秒。

换成无量纲坐标后，方程中的黏度系数变为$1/\mathrm{Re}$，其中$\mathrm{Re}=U_0L_0/\nu$是雷诺数。它由所选尺度和黏度共同确定。实现时可直接对无量纲坐标求导，使用无量纲方程计算残差。

本章的二维模型用于学习基本流程。三维重建时，网络输入扩展为$(x,y,z,t)$，速度增加$w$分量，连续性方程相应增加$\partial_zw$。若研究时间平均风场，还需要考虑湍流脉动对平均流动的影响，常通过湍流模型来描述。后续深入流体建模时再展开这一部分。

\section{把数据和方程放进同一次训练}
\label{sec:pinn-training}

\subsection{两条计算路径汇入一个损失}
图\ref{fig:pinn-constraint-path}展示了风场 PINN 的主要计算过程。网络先根据坐标给出速度和压力，然后分别计算径向速度误差与方程残差，最后合成损失来更新同一组网络参数。

\begin{figure}[htbp]
\centering
\begin{tikzpicture}[font=\small]
 \node[idi input, minimum width=2.2cm] (coord) at (0,0) {时空坐标 $(x,y,t)$};
 \node[idi process, minimum width=2.2cm] (net) at (0,-1.1) {网络 $F_\theta$};
 \node[idi state, minimum width=3.0cm] (field) at (0,-2.2) {预测 $u_\theta,v_\theta,\pi_\theta$};
 \node[idi process, minimum width=2.7cm] (proj) at (-3,-3.6) {观测算子\\预测径向速度};
 \node[idi process, minimum width=2.7cm] (pde) at (3,-3.6) {自动微分\\方程残差};
 \node[idi input, minimum width=1.7cm] (obs) at (-5,-5) {雷达读数};
 \node[idi output, minimum width=2.2cm] (data) at (-2,-5) {数据损失};
 \node[idi output, minimum width=2.2cm] (physics) at (3,-5) {物理损失};
 \node[idi output, minimum width=2.8cm] (loss) at (0,-6.3) {总损失 $\mathcal L$};
 \draw[idi arrow] (coord)--(net);
 \draw[idi arrow] (net)--(field);
 \draw[idi arrow] (field)--(proj);
 \draw[idi arrow] (field)--(pde);
 \draw[idi arrow] (proj)--(data);
 \draw[idi arrow] (obs)--(data);
 \draw[idi arrow] (pde)--(physics);
 \draw[idi arrow] (data)--(loss);
 \draw[idi arrow] (physics)--(loss);
 \draw[idi feedback] (loss.east)--(5.2,-6.3)--(5.2,-1.1)--node[above]{更新参数} (net.east);
\end{tikzpicture}
\caption{风场 PINN 的计算路径。实线表示预测量和误差的计算，虚线表示根据总损失更新网络参数。可用的初值、边界信息也可形成损失项，一起加入总损失。}
\label{fig:pinn-constraint-path}
\end{figure}

为了计算二阶导数，网络通常选用平滑激活函数，如$\tanh$。第一次实现时，可以先从普通全连接网络开始，把观测投影和方程残差这两条路径分别检查清楚。

\subsection{数据损失检查预测与测量是否接近}
在第$m$条观测的位置查询网络，用观测算子计算$\widehat y_m$，再与实测径向速度$y_m$比较。若以$U_0$缩放速度，数据损失可写成
\begin{equation}
 \mathcal L_d=\frac1{N_d}\sum_{m=1}^{N_d}
 \left(\frac{\widehat y_m-y_m}{U_0}\right)^2.
 \label{eq:pinn-weighted-data-loss}
\end{equation}
$N_d$是观测条数，平均平方误差越小，预测读数与实际读数越接近。例如实测值是$3\,\mathrm{m/s}$、预测值是$2.8\,\mathrm{m/s}$，两者先在相同单位下相减，再除以速度尺度。

\subsection{物理损失检查预测是否符合方程}
用$\overline r_c,\overline r_u,\overline r_v$表示无量纲方程残差。在$N_f$个配点上计算
\begin{equation}
 \mathcal L_f=\frac1{N_f}\sum_{j=1}^{N_f}
 \left[\overline r_c(\boldsymbol\xi_j)^2+
 \overline r_u(\boldsymbol\xi_j)^2+
 \overline r_v(\boldsymbol\xi_j)^2\right],
\end{equation}
其中$\boldsymbol\xi_j=(x_j^*,y_j^*,t_j^*)$是一个配点。这里与手算例子的做法相同，只是从一条方程扩展到了连续性和两个动量方程。

如果从有量纲残差出发，也可以用$\overline r_c=(L_0/U_0)r_c$、$\overline r_u=(L_0/U_0^2)r_u$和$\overline r_v=(L_0/U_0^2)r_v$得到相同的尺度处理。

若已知边界或初始状态，还可以计算边界损失$\mathcal L_b$和初值损失$\mathcal L_i$。一个常见的总损失为
\begin{equation}
 \mathcal L=\lambda_d\mathcal L_d+\lambda_f\mathcal L_f
 +\lambda_b\mathcal L_b+\lambda_i\mathcal L_i.
 \label{eq:pinn-loss}
\end{equation}
各个$\lambda$是权重，用来调整训练对不同误差的重视程度。先把各项缩放到合适的量级，再观察数据误差和方程残差怎样下降，会更容易判断权重是否合适。

\subsection{一次训练和一张风场图怎样得到}
训练通常从随机网络参数开始。每一步抽取观测点和配点，分别计算两类误差，用 Adam 等优化器更新参数，然后重复。训练过程中分别查看各项损失，能帮助定位投影、求导或尺度处理中的问题。

\begin{algorithm}[htbp]
\caption{从稀疏雷达观测训练风场 PINN}
\label{alg:pinn-wind}
\begin{algorithmic}[1]
\Require 测点坐标、时间、波束方向、径向速度，以及选定的流体方程
\State 统一坐标与单位，对坐标和速度做无量纲化
\State 建立网络，在重建区域内准备配点
\For{每一步训练}
 \State 从观测中取一批样本，计算预测的径向速度和数据损失
 \State 从配点中取一批坐标，自动微分计算方程残差和物理损失
 \State 加入可用的边界、初值误差，计算总损失
 \State 根据总损失更新网络参数
\EndFor
\State 固定一个查询时刻，在空间网格上查询网络
\State 恢复物理单位，绘制速度分量、风速大小或风向
\end{algorithmic}
\end{algorithm}

例如想画$t=50\,\mathrm{s}$时的风场，就把这个时刻与网格上的各个$(x_i,y_i)$组合起来，逐点或成批送入网络。得到$u_i,v_i$后，计算风速大小$V_i=\sqrt{u_i^2+v_i^2}$，用颜色展示，同时用箭头表示速度方向。这样就从一组稀疏读数得到了可查看的空间分布。

检查效果时，既看训练误差，也与未参与拟合或模型选择的观测、参考风场比较。训练数据损失只回答已用雷达读数是否拟合得好。参考风场对比可检查未测区域和速度分量的重建误差。若没有独立参考场，低方程残差和低训练误差仍不能证明真实风场已被唯一恢复。


\section{双束激光雷达风场重建案例}
\label{sec:pinn-validation}

\subsection{先看任务、测点与参考风场}
Zhang 和 Zhao 在 2021 年的研究中，将 PINN 用于风机前方的时空风场重建\cite{zhang2021spatiotemporal}。他们关注轮毂高度处的水平面，希望从两束激光雷达的径向速度，得到整个平面内的顺风速度$u$和横风速度$v$。

图\ref{fig:pinn-zhang-measurements}摘录了论文的测量布置。图 A 中，安装在风机上的雷达向上游发出两束激光，夹角为$30^\circ$，每束相对中心方向偏转$15^\circ$。叉号表示测点，蓝色区域是希望恢复风场的平面。图 B 是左右波束在不同测点得到的径向风速时间序列。训练数据是这些分散位置上的曲线，重建目标是整片区域随时间变化的速度分布。

\begin{figure}[htbp]
\centering
\includegraphics[width=\linewidth]{images/pinn-zhang2021/fig1-measurements.pdf}
\caption{双波束激光雷达的测量布置与径向速度记录。A 为测量几何，B 为各测点的时间序列。摘自 Zhang 和 Zhao\cite{zhang2021spatiotemporal}原文图 1(A、B)，保留原图标注。}
\label{fig:pinn-zhang-measurements}
\end{figure}

研究使用 SOWFA 大涡模拟生成风场：先模拟三维大气边界层，再取轮毂高度处的二维切片作为参考。从两束波束经过的位置提取速度，投影后得到模拟雷达读数。训练时使用这些径向读数，完整切片用于对比重建结果。这种数据组织使我们既有稀疏观测，也有可逐点比较的参考风场。

论文中用于重建的区域为$x\in[-240,0]\,\mathrm{m}$、$y\in[-60,60]\,\mathrm{m}$，时间窗为$0$到$100\,\mathrm{s}$。表\ref{tab:pinn-zhang-setup}列出了基准算例的主要设置。

\begin{table}[htbp]
\centering\small
\caption{论文基准算例的主要设置，整理自原文第 2.3、3.2 节与表 2\cite{zhang2021spatiotemporal}。}
\label{tab:pinn-zhang-setup}
\begin{tabular}{m{0.33\linewidth}m{0.58\linewidth}}
\toprule
项目 & 设置\\
\midrule
平均来流与湍流强度 & $8\,\mathrm{m/s}$，自由来流湍流强度 $6\%$\\
雷达波束 & 两束，半夹角 $15^\circ$，测量范围 $220\,\mathrm{m}$\\
空间与时间采样 & 每束 11 个测点，相邻测点相距 $20\,\mathrm{m}$，每秒测量一次\\
基准测量噪声 & 不添加噪声\\
网络 & 含输入、输出在内共 12 层，隐藏层宽度 128，$\tanh$ 激活\\
优化器与学习率 & Adam，$10^{-4}$\\
每步训练样本 & 物理配点 1000 个，左右波束观测各 1100 条\\
\bottomrule
\end{tabular}
\end{table}

这里的“1100 条观测”统计的是位置与时间的组合。每束只有 11 个空间测点，但连续测量后，每个测点都会形成一条时间序列。论文还在$81\times41\times101$的时空网格上准备配点坐标，每一步从中抽取一批来计算方程残差。

\subsection{网络为什么先输出流函数和压力}
这个案例的基础网络输出\textbf{流函数}$\psi$和压力，再通过求导得到速度。为沿用本章符号，将论文的无量纲压力$p$记为$\pi$，把网络写成
\begin{equation}
 (\psi_\theta,\pi_\theta)=F_\theta(x,y,t).
 \label{eq:pinn-zhang-network}
\end{equation}
本小节及后续案例公式中的坐标、时间和输出均已无量纲化。论文选择风机直径$D$作为长度尺度、平均来流速度$U_\infty$作为速度尺度，时间尺度就是$D/U_\infty$。

流函数可以看成一个用于生成速度的辅助函数。按照论文原文式 (3)，由它计算
\begin{equation}
 u_\theta=\frac{\partial\psi_\theta}{\partial y},\qquad
 v_\theta=-\frac{\partial\psi_\theta}{\partial x}.
 \label{eq:pinn-zhang-streamfunction}
\end{equation}
这样做的好处是，连续性方程会自动满足：
\[
 \partial_xu_\theta+\partial_yv_\theta
 =\partial_x\partial_y\psi_\theta-\partial_y\partial_x\psi_\theta=0.
\]
例如取$\psi=2y$，就得到均匀速度$u=2,v=0$。再把$\psi$换成可以学习的复杂函数，就能表示随空间、时间变化的速度分布，同时保留上述连续性关系。

图\ref{fig:pinn-zhang-network}摘录了论文的网络结构。灰色 DNN1 表示基础网络，输出$\psi,p$。绿色 DNN2 通过自动微分得到$u,v$，再通过$f_1,f_2$得到左右波束的径向速度。粉色 DNN3 进一步计算动量方程残差。图中的三个 DNN 名称对应同一组可训练参数在不同计算步骤中的用法，训练时一起更新。

\begin{figure}[H]
\centering
\includegraphics[width=\linewidth]{images/pinn-zhang2021/fig1-network.pdf}
\caption{论文中的网络及损失计算。$\mathcal L_1$来自方程残差，$\mathcal L_2$来自径向观测误差，两者共同训练基础网络。原图用$p$表示无量纲压力，对应本章的$\pi$。摘自 Zhang 和 Zhao\cite{zhang2021spatiotemporal}原文图 1(C)。}
\label{fig:pinn-zhang-network}
\end{figure}

\subsection{两束观测怎样参与损失计算}
网络得到$u,v$后，先按波束方向计算径向速度。沿用论文的左右波束与符号约定，原文式 (5)--(6) 为
\begin{align}
 f_r(u,v)&=u\cos15^\circ-v\sin15^\circ,\\
 f_l(u,v)&=u\cos(-15^\circ)-v\sin(-15^\circ).
 \label{eq:pinn-zhang-projection}
\end{align}
右波束使用$f_r$，左波束使用$f_l$。例如在右波束某个位置查询得到$u=1,v=0.1$，预测径向速度约为$1\times0.966-0.1\times0.259=0.940$，再与该位置的无量纲测量值比较。

记右、左波束实测值分别为$y_{r,i}$和$y_{l,i}$，对应预测值为$\widehat y_{r,i}$和$\widehat y_{l,i}$。论文的数据损失是左右两组平均平方误差相加：
\begin{equation}
 \mathcal L_2=\frac1{N_r}\sum_{i=1}^{N_r}(\widehat y_{r,i}-y_{r,i})^2
 +\frac1{N_l}\sum_{i=1}^{N_l}(\widehat y_{l,i}-y_{l,i})^2.
 \label{eq:pinn-zhang-data}
\end{equation}
$N_r,N_l$是本次训练取出的两组观测条数。每个预测值都在其对应的测点位置和测量时刻计算。

物理部分使用二维无量纲动量方程。按照原文式 (9)--(10)，两个残差为
\begin{align}
 e_u&=\partial_tu+u\partial_xu+v\partial_yu+\partial_x\pi
 -\frac1{\mathrm{Re}}(\partial_{xx}u+\partial_{yy}u),\\
 e_v&=\partial_tv+u\partial_xv+v\partial_yv+\partial_y\pi
 -\frac1{\mathrm{Re}}(\partial_{xx}v+\partial_{yy}v).
 \label{eq:pinn-zhang-residual}
\end{align}
其中$\mathrm{Re}=U_\infty D/\nu$，$U_\infty$是参考来流速度，$D$是参考长度，$\nu$是运动黏度。在$N_{\mathrm{NS}}$个配点上取残差平方的平均，得到
\begin{equation}
 \mathcal L_1=\frac1{N_{\mathrm{NS}}}\sum_{j=1}^{N_{\mathrm{NS}}}
 \big[e_u(\boldsymbol\xi_j)^2+e_v(\boldsymbol\xi_j)^2\big].
 \label{eq:pinn-zhang-physics}
\end{equation}
论文最后采用$\mathcal L=\mathcal L_1+\mathcal L_2$训练网络。流函数构造在速度场足够光滑时满足连续性，所以物理损失中只列出两个动量残差。该案例明确列出的训练目标由径向数据项与域内动量残差项组成，没有另外列出初值损失和边界损失。这一设置依赖数据覆盖和所选方程，不能视为显式施加了边界条件。

把这一过程连起来看：一批坐标经过网络得到$\psi,\pi$，求导得到$u,v$。接着，在观测点计算径向误差，在配点计算动量残差。将两类误差相加，再更新同一组参数。训练使用的风速标签来自两束波束，压力和流函数由训练过程一同确定。

\subsection{从风场图片看重建效果}
训练完成后，固定时刻，在$81\times41$个空间网格点上查询$u,v$，即可得到一张风场图。图\ref{fig:pinn-zhang-baseline}给出基准算例的结果：左列是网络重建，右列是 SOWFA 参考切片，三行分别对应$50$、$60$和$70\,\mathrm{s}$。

\begin{figure}[H]
\centering
\includegraphics[width=0.92\linewidth]{images/pinn-zhang2021/fig3-baseline.pdf}
\caption{基准算例的重建与参考风场。左列为 PINN 结果，右列为 SOWFA 参考值，从上到下为$50$、$60$、$70\,\mathrm{s}$。颜色表示风速大小，箭头表示速度方向，叉号表示雷达测点。摘自 Zhang 和 Zhao\cite{zhang2021spatiotemporal}原文图 3。}
\label{fig:pinn-zhang-baseline}
\end{figure}

看这张图时，可以先找红黄色的高速区域。它在$50\,\mathrm{s}$时位于区域左侧，$60\,\mathrm{s}$时移到中部，$70\,\mathrm{s}$时到达风机附近。左列与右列呈现出相近的移动过程，说明重建结果抓住了这团高速来流的位置和输运过程。再看箭头，可以比较两个结果中的局部风向。

除了看图，论文还逐时刻计算全场风速误差，再对时间取平均。将风速大小记为$V=\sqrt{u^2+v^2}$，对应原文式 (14) 的指标为
\begin{equation}
 \varepsilon_V=\frac1{N_t}\sum_{k=1}^{N_t}
 \sqrt{\frac1{N_s}\sum_{i=1}^{N_s}
 \left(\widehat V_{i,k}-V^{\mathrm{ref}}_{i,k}\right)^2}.
 \label{eq:pinn-zhang-mrmse}
\end{equation}
先在一个时刻，对$N_s$个空间位置求均方根误差，再把$N_t$个时刻的结果取平均。这就是论文中的 MRMSE。本例$N_s=81\times41=3321$，指标采用 100 个时刻。风向使用同样的计算顺序。

基准算例的风速大小 MRMSE 为$0.198\,\mathrm{m/s}$，风向 MRMSE 为$2.77^\circ$。表\ref{tab:pinn-zhang-results}还整理了几个改变观测条件的结果，便于读者看到测量质量和测点数量怎样影响重建。

\begin{table}[htbp]
\centering\small
\caption{部分算例的全场重建误差，整理自 Zhang 和 Zhao\cite{zhang2021spatiotemporal}原文表 3。各行采用同一 MRMSE 计算方式。}
\label{tab:pinn-zhang-results}
\begin{tabular}{m{0.46\linewidth}rr}
\toprule
测量或来流设置 & 风速误差（$\mathrm{m/s}$） & 风向误差（$^\circ$）\\
\midrule
基准：无噪声，每束 11 个测点 & 0.198 & 2.77\\
径向速度加均匀噪声$[-0.1,0.1]\,\mathrm{m/s}$ & 0.387 & 3.73\\
空间采样变稀：每束 6 个测点，间隔 $40\,\mathrm{m}$ & 0.212 & 2.85\\
时间采样变稀：每 $2\,\mathrm{s}$ 测量一次 & 0.222 & 2.66\\
来流相对雷达中心方向偏转 $20^\circ$ & 0.281 & 2.46\\
\bottomrule
\end{tabular}
\end{table}

加入噪声后，风速误差有所增加。减少空间测点后，模型仍重建出了主要风场结构。学习这个案例时，可以把注意力放在这种联系上：输入数据发生了什么变化，重建图中的结构和误差数值又怎样变化。

\subsection{从重建时间窗走向短时预测}
前面的对比发生在$0$到$100\,\mathrm{s}$的测量时间窗内。论文还把查询时间延伸到$115\,\mathrm{s}$，用额外 15 秒的 SOWFA 结果检验短时预测。

图\ref{fig:pinn-zhang-timeseries}给出风机位置附近三个横向位置的风速曲线。黑色实线为参考值，红色虚线为网络预测，横轴延伸到$115\,\mathrm{s}$。读图时可分别观察$100\,\mathrm{s}$之前的重建段，以及之后的短时外推段，比较峰值、低谷和变化趋势的对应情况。

\begin{figure}[htbp]
\centering
\includegraphics[width=\linewidth]{images/pinn-zhang2021/fig5-time-series.pdf}
\caption{三个横向位置的风速时间序列。横向坐标依次为$y=0$、$15$、$30\,\mathrm{m}$。黑色实线为参考值，红色虚线为预测值。曲线包含$100$到$115\,\mathrm{s}$的未来时间外推。摘自 Zhang 和 Zhao\cite{zhang2021spatiotemporal}原文图 5。}
\label{fig:pinn-zhang-timeseries}
\end{figure}

这个案例把 PINN 的基本想法落到了具体计算上：用坐标网络表示风场，用流函数构造满足连续性方程的速度分布，用投影连接雷达读数，用动量残差联系未测位置，最后在网格和时间轴上查询结果。回到自己的研究时，可以沿着同一顺序整理测量几何、方程、训练数据和结果对比。

\FloatBarrier
\section{PINN 的优化方向}
\label{sec:pinn-optimization}
前面的案例给出了一个完整的 PINN：用网络表示风场，再用观测误差和方程残差训练它。真正动手后，还会遇到一些具体问题。例如，数据误差下降很快，方程残差却迟迟不降。有时平均风速已经接近参考值，局部变化仍显得过于平滑。也可能短时间内重建较好，时间窗一拉长，结果就变差。本节围绕这些现象，介绍怎样调整训练方法和网络结构。

Krishnapriyan 等的研究帮助我们理解训练困难从哪里来\cite{krishnapriyan2021failure}。Wang 等整理了尺度处理、网络设计和训练策略方面的实用方法\cite{wang2023expert}。Jagtap 和 Karniadakis 提出的 XPINNs 则把一个大问题拆成多个相互连接的子问题\cite{jagtap2020xpinn}。下面先从容易检查的设置开始，再介绍更深入的改进。

\subsection{先看清训练卡在哪里}
一个网络要同时拟合雷达读数、满足连续性和动量方程，还可能需要匹配入口与初始状态。这几项任务使用同一组参数，但各自提出的调整要求会有所不同。因此，查看训练过程时，可以先把总损失拆开看。

例如，若$\mathcal L_d$持续下降而$\mathcal L_f$几乎不变，就需要检查方程残差的尺度和权重。若两项都下降，而高速区域仍被预测得很模糊，可以继续检查配点分布与网络表达细节的能力。把问题定位到某一类现象后，再选择对应方法，调整会更有针对性。

Krishnapriyan 等在对流、反应和反应扩散方程上发现，增大方程中的某些系数后，原先能够训练好的 PINN 会变得难以优化。其原因与损失函数的形状有关：可以把参数训练想成在起伏的地形上寻找低点，方程项会改变这片地形，使下降过程更难进行。论文通过改变训练方式，在相同网络结构下得到了更好的结果\cite{krishnapriyan2021failure}。这为我们提供了一个很实际的起点：先检查训练安排，再决定是否增加网络规模。

\subsection{统一变量尺度，并把配点放到合适的位置}
\paragraph{先处理尺度。}
第\ref{sec:pinn-physics}节已经介绍了无量纲化。它在优化中的作用，可以理解为让网络在相近的数值范围内处理坐标、时间、速度和压力。例如，用$L_0=240\,\mathrm{m}$、$U_0=8\,\mathrm{m/s}$作为一组教学尺度，则$x=-120\,\mathrm{m}$对应$x^*=-0.5$，$t=30\,\mathrm{s}$对应$t^*=1$，$u=8\,\mathrm{m/s}$对应$u^*=1$。这些量更便于放进同一个网络。

实际操作时，要把同一套尺度同时用于观测投影、方程残差和结果恢复。如果另外把网络输入坐标映射到$[-1,1]$，求导计算也应包含这次坐标变换。Wang 等把无量纲化放在训练流程的第一步，并在 Stokes 流动算例中比较了它对训练结果的影响\cite{wang2023expert}。

\paragraph{让配点覆盖整个重建区域。}
配点决定了网络在哪些位置接受方程检查。对于稀疏雷达重建，观测集中在几条波束上，物理配点则可以覆盖波束之间和波束之外的区域。这样，方程信息就能进入整个待重建区域。

一种容易实现的做法是每一步随机抽取一批配点，或每隔若干步重新采样，使网络不断在新的位置接受检查。Wang 等将随机采样纳入了训练流程\cite{wang2023expert}。如果局部剪切较强，还可以进一步使用\textbf{残差自适应采样}：先在一批候选点上计算残差，找出较大的位置，再给这些区域增加配点。保留一部分均匀分布的点，可以同时照顾整体区域与局部细节。

在风场问题中，这一策略可以按“全域采样、查看残差分布、局部加密、继续训练”的顺序实施。增加的是方程检查位置，所需输入仍然只是坐标和时间。

\subsection{平衡数据损失和物理损失}
即使变量已经无量纲化，不同损失对参数更新的影响仍可能相差很大。这里要关注的是\textbf{损失对参数的梯度}，也就是每项损失在多大程度上推动网络调整。某项损失的梯度长期占主导时，训练就会更多地围绕这项任务进行。

Wang 等介绍了一种基于梯度大小的自适应加权方法\cite{wang2023expert}。用$q$表示数据、物理、边界或初值等某一项损失，先计算
\begin{equation}
 g_q=\left\|\nabla_\theta\mathcal L_q\right\|_2,\qquad
 \widehat\lambda_q=\frac{\sum_j g_j}{g_q}.
 \label{eq:pinn-gradient-balance}
\end{equation}
$\nabla_\theta\mathcal L_q$是该项损失对全部网络参数的梯度，$g_q$表示它的整体大小。梯度较小的项会获得较大的权重，使各项对参数更新的推动幅度更接近。实际实现需对很小的$g_q$加入数值稳定项或裁剪权重，否则当$g_q=0$时该定义无效，极小梯度也会产生过大的权重。

例如只有数据和物理两项时，若$g_d=0.01$、$g_f=1$，按上式得到$\widehat\lambda_d=101$、$\widehat\lambda_f=1.01$。加权后的梯度大小均为$1.01$。这个例子说明，权重是在协调各项任务参与训练的程度。

为使权重变化平滑，可以采用移动平均：
\begin{equation}
 \lambda_q^{\mathrm{new}}
 =\alpha\lambda_q^{\mathrm{old}}+(1-\alpha)\widehat\lambda_q.
 \label{eq:pinn-gradient-balance-ema}
\end{equation}
$\alpha$越大，越多地保留原来的权重。Wang 等给出$\alpha=0.9$、每 1000 步更新一次权重的默认设置\cite{wang2023expert}。权重定期重算，在当前一步更新网络参数时，将它们作为固定系数使用。

对于风场重建，可以先分别记录$\mathcal L_d$、连续性残差和两个动量残差。若某一项长期下降缓慢，再比较梯度大小，决定按“数据与物理”两组加权，还是把不同方程进一步分开加权。这样，调节权重就有了可观察的依据。

\subsection{采用课程训练，让任务由易到难}
课程训练的想法很接近学习过程：先完成较容易的任务，再利用已经学到的参数处理更难的任务。每个阶段都以上一阶段的网络为起点，逐步接近最终问题。

Krishnapriyan 等以一维对流方程为例：
\begin{equation}
 \partial_tq+\beta\partial_xq=0.
 \label{eq:pinn-curriculum-convection}
\end{equation}
$q$是随时间平移的一个场，$\beta$控制它的平移速度。论文先在较小的$\beta$下训练，再逐步增大到目标值。原文表 1 中，$\beta=30$时，常规训练的相对误差为$8.97\times10^{-1}$，课程训练为$2.02\times10^{-2}$\cite{krishnapriyan2021failure}。这个例子直观地说明，训练顺序本身就能影响最后学到的结果。

Wang 等也介绍了类似做法：求解较高雷诺数的流动时，先训练较低雷诺数的模型，再把结果用于目标雷诺数下的训练\cite{wang2023expert}。前一阶段提供一个较好的起点，后一阶段继续拟合目标方程。

把“由易到难”的思路用于风场重建，可以设计分阶段训练：先让网络学到主要来流和观测变化，再逐步增强动量方程约束，最后在目标权重和全部观测上继续训练。如果采用改变方程参数的课程，各阶段也要明确当前使用的参数，最后回到目标方程。这样就能把每一阶段的训练任务说清楚。

\subsection{缩短单次时间跨度，并按时间顺序学习}
\paragraph{时间分窗：每次重建一小段。}
假设有$0$到$100\,\mathrm{s}$的雷达数据，可以把它划成五个$20\,\mathrm{s}$窗口。先训练第一个窗口，得到$t=20\,\mathrm{s}$的场，再把这一末态作为下一个窗口的初始状态，同时使用下一窗口内的雷达观测继续重建。网络参数也可以沿用前一窗口的训练结果。

这就是 Krishnapriyan 等讨论的序列式学习思路：将一次学习整个时空区域，改成沿时间逐段推进\cite{krishnapriyan2021failure}。Wang 等也将时间推进用于长时间动力系统的训练\cite{wang2023expert}。对风场而言，每个网络只需处理一段较短时间内的变化，窗口之间再通过末态与初值连接。

\paragraph{因果加权：在一个窗口内先学好较早的时刻。}
时间分窗以后，每个窗口内部仍包含许多时刻。因果加权进一步调整这些时刻的学习顺序。较早时间段的残差还很大时，先把训练重点放在前面。等前面的残差下降后，再逐步提高后面时间段的权重。

按照 Wang 等的写法，将当前窗口划成$M$个有先后顺序的小段，用$\mathcal L_f^{(k)}$表示第$k$段的物理损失，定义
\begin{align}
 w_1&=1,\qquad
 w_k=\exp\left(-\varepsilon\sum_{j=1}^{k-1}\mathcal L_f^{(j)}\right),
 \quad k=2,\ldots,M,\\
 \mathcal L_{f,\mathrm{causal}}
 &=\frac1M\sum_{k=1}^{M}w_k\mathcal L_f^{(k)}.
 \label{eq:pinn-causal-weight}
\end{align}
$\varepsilon>0$控制这种先后顺序的强弱。前面各段的累计残差越大，后面段的$w_k$越小。前面的残差下降后，后面的权重就逐渐接近$1$。更新网络参数时，同样把当前计算出的$w_k$作为固定系数\cite{wang2023expert}。

例如取$\varepsilon=1$，若前两段物理损失都是$1$，三段权重约为$(1,0.37,0.14)$。当前两段损失都降到$0.1$时，权重变为$(1,0.90,0.82)$。这些教学数值展示了训练重点怎样逐渐向后展开。

图\ref{fig:pinn-time-strategies}把两种做法放在一起。上面是不同时间窗口之间传递末态，下面是同一个网络内部按时间段调整物理损失权重。它们可以结合使用：先把长记录分窗，再在每个窗口内进行因果加权。

\begin{figure}[htbp]
\centering
\begin{tikzpicture}[font=\small]
 \node at (4,0.9) {时间分窗：每个窗口分别训练};
 \node[idi process, minimum width=3.0cm] (win1) at (0,0) {$F_{\theta_1}$\\$0$--$20\,\mathrm{s}$};
 \node[idi process, minimum width=3.0cm] (win2) at (4,0) {$F_{\theta_2}$\\$20$--$40\,\mathrm{s}$};
 \node[idi process, minimum width=3.0cm] (win3) at (8,0) {$F_{\theta_3}$\\$40$--$60\,\mathrm{s}$};
 \draw[idi arrow] (win1)--node[above,font=\footnotesize]{末态}(win2);
 \draw[idi arrow] (win2)--node[above,font=\footnotesize]{末态}(win3);
 \node at (4,-1.15) {因果加权：同一窗口、同一网络内部};
 \node[idi state, minimum width=3.0cm] (seg1) at (0,-2.05) {$I_1$：较早时间段\\权重 $w_1=1$};
 \node[idi state, minimum width=3.0cm] (seg2) at (4,-2.05) {$I_2$：中间时间段\\权重 $w_2$};
 \node[idi state, minimum width=3.0cm] (seg3) at (8,-2.05) {$I_3$：较晚时间段\\权重 $w_3$};
 \draw[idi arrow] (seg1)--node[above,font=\footnotesize]{影响}(seg2);
 \draw[idi arrow] (seg2)--node[above,font=\footnotesize]{影响}(seg3);
 \node[align=center,font=\footnotesize] at (4,-3.05) {较早时间段的累计残差下降后，后面时间段的权重逐步增大};
\end{tikzpicture}
\caption{时间分窗与因果加权的教学示意。上排箭头表示前一窗口的末态用于下一窗口的初值。下排箭头表示较早时间段的残差参与后面权重的计算，见式\eqref{eq:pinn-causal-weight}。图中窗口长度用于举例。}
\label{fig:pinn-time-strategies}
\end{figure}

\subsection{改进网络对局部细节的表达}
风场往往同时包含缓慢变化的背景来流和变化较快的局部结构。普通全连接网络在训练中倾向于先学到平滑、变化缓慢的部分，这种现象称为\textbf{频谱偏置}。它可以帮助我们理解一种常见结果：平均风速已经比较接近，局部峰值和剪切变化却仍被抹平。Wang 等从以下几方面改进网络\cite{wang2023expert}。

\paragraph{傅里叶特征：给坐标增加不同频率的表示。}
在坐标进入网络前，先把它转换为一组正弦和余弦特征：
\begin{equation}
 \gamma(\boldsymbol\xi)=
 \begin{bmatrix}
  \cos(B\boldsymbol\xi)\\
  \sin(B\boldsymbol\xi)
 \end{bmatrix}.
 \label{eq:pinn-fourier-features}
\end{equation}
$\boldsymbol\xi$是无量纲时空坐标，$B$是一组控制频率的系数。随机傅里叶特征通常从高斯分布采样这些系数，再在训练中固定它们。以一个空间坐标为例，$\sin(2\pi x^*)$在单位区间内变化一周，$\sin(10\pi x^*)$则变化五周。将这些特征交给网络，相当于提供不同快慢的变化模式，让它更容易组合出局部细节。

频率范围需要与希望重建的空间、时间尺度相配合。在风场中，可以比较加入特征前后，高速区边缘、横向剪切和速度时间曲线是否更清晰。这里增加的是坐标的表示方式，网络的输出仍然是本章使用的速度、压力或流函数。

\paragraph{随机权重分解：调整参数的组织方式。}
Wang 等还介绍了随机权重分解（RWF）。对于一层网络的权重矩阵，可以写成
\begin{equation}
 W=\operatorname{diag}(\exp(\mathbf a))V.
 \label{eq:pinn-rwf}
\end{equation}
$\mathbf a$为每个神经元提供一个可学习的缩放量，$V$给出其余权重。先按普通方法初始化$W$，再选取$\mathbf a$并据此得到$V$，使分解后的初始矩阵与原来的$W$一致。训练时更新$\mathbf a,V$。这种改法保留了全连接层的计算形式，同时改变参数的更新过程，使不同神经元可以形成不同的调整幅度。

\paragraph{改进连接方式，并利用已知物理结构。}
Wang 等讨论的 modified MLP 用两个编码器提取输入特征，再在各隐藏层中组合这些特征，使输入信息能反复参与中间层计算。它提供了普通全连接网络之外的另一种结构选择\cite{wang2023expert}。

物理关系也可以直接参与输出设计。第\ref{sec:pinn-validation}节的流函数就是一个例子：网络输出$\psi$，再由求导得到$u,v$，从而把连续性方程放进速度的构造过程。若问题具有已知的周期边界，还可以通过周期性的正弦、余弦输入来表达这种边界关系。这样的结构设计能减少需要单独学习的约束。

\subsection{用 XPINNs 把大区域拆成相互连接的子区域}
当重建区域变大，或区域内不同位置的变化程度差别明显时，可以让多个局部网络分工。Jagtap 和 Karniadakis 提出的 XPINNs 将空间或时空区域划成若干子区域，每个区域使用一个网络，再通过界面条件把结果连接起来\cite{jagtap2020xpinn}。

图\ref{fig:pinn-xpinn-domain}沿用前面案例的矩形区域，示意在$x=-120\,\mathrm{m}$处划分。左侧由$F_{\theta_1}$表示，右侧由$F_{\theta_2}$表示。各自使用区域内的观测和配点训练，共同边界$\Gamma$上的点则同时送入两个网络，比较两侧的结果。

\begin{figure}[htbp]
\centering
\begin{tikzpicture}[font=\small]
 \fill[blue!7] (-5.2,-1.25) rectangle (0,1.25);
 \fill[orange!10] (0,-1.25) rectangle (5.2,1.25);
 \draw[thick] (-5.2,-1.25) rectangle (5.2,1.25);
 \draw[thick,dashed] (0,-1.25)--(0,1.25);
 \foreach \yy in {-1.0,-0.5,0,0.5,1.0} {\fill (0,\yy) circle (1.6pt);}
 \node[align=center] at (-2.6,0) {子区域 $\Omega_1$\\网络 $F_{\theta_1}$};
 \node[align=center] at (2.6,0) {子区域 $\Omega_2$\\网络 $F_{\theta_2}$};
 \draw[idi arrow] (-4,1.8)--node[above]{来流方向}(4,1.8);
 \node[below] at (-5.2,-1.3) {$x=-240\,\mathrm{m}$};
 \node[below] at (0,-1.3) {$-120\,\mathrm{m}$};
 \node[below] at (5.2,-1.3) {$0\,\mathrm{m}$};
 \node[align=center,font=\footnotesize] at (0,-2.35) {界面 $\Gamma$ 上的黑点：在同一坐标、同一时刻\\比较两个网络的预测与方程残差};
\end{tikzpicture}
\caption{风场区域分解的教学示意。虚线为两个子区域的界面，黑点为界面约束的计算位置，箭头表示来流方向。这里的划分用于解释 XPINNs 的计算方法。}
\label{fig:pinn-xpinn-domain}
\end{figure}

界面上主要处理两件事：\textbf{两侧给出相衔接的场值，两侧的方程残差也保持一致}。沿用本章的无量纲变量，记第$k$个局部网络的预测为$\mathbf z_k=(u_k,v_k,\pi_k)^\top$，残差为$\mathbf R_k=(\overline r_{c,k},\overline r_{u,k},\overline r_{v,k})^\top$。

先在界面点上计算两侧预测的平均值
\[
 \overline{\mathbf z}=\frac{\mathbf z_1+\mathbf z_2}{2}.
\]
按 XPINNs 的界面构造思路，让每一侧接近这个共同值：
\begin{equation}
 \mathcal L_{\Gamma,z}=\frac1{N_\Gamma}\sum_{i=1}^{N_\Gamma}
 \left[
 \|\mathbf z_1(\boldsymbol\xi_i)-\overline{\mathbf z}(\boldsymbol\xi_i)\|_2^2
 +\|\mathbf z_2(\boldsymbol\xi_i)-\overline{\mathbf z}(\boldsymbol\xi_i)\|_2^2
 \right].
 \label{eq:pinn-xpinn-interface-value}
\end{equation}
同时，比较同一界面点上的残差：
\begin{equation}
 \mathcal L_{\Gamma,R}=\frac1{N_\Gamma}\sum_{i=1}^{N_\Gamma}
 \|\mathbf R_1(\boldsymbol\xi_i)-\mathbf R_2(\boldsymbol\xi_i)\|_2^2.
 \label{eq:pinn-xpinn-interface-residual}
\end{equation}
$N_\Gamma$是界面点数，$\boldsymbol\xi_i$记录其坐标和时间。第一项帮助两块预测衔接，第二项协调界面处的方程计算。它们与各子区域原有的数据损失、物理损失一起参与训练。具体方程需要时，还可以加入通量等界面条件\cite{jagtap2020xpinn}。

XPINNs 的另一个特点是不同区域可以采用不同配置。变化平缓的区域可以用较小的网络，局部变化较强的区域可以增加网络容量和配点。多个子网络还可以分配到不同计算设备上训练，这样计算资源就能按区域需要分配。

空间分解与时间分窗也有不同侧重：前者让不同空间区域的网络通过界面共同训练，后者沿时间把前一窗口的末态传给后一窗口。XPINNs 的划分可以同时涉及空间和时间，因此能够容纳更灵活的时空子区域。

\subsection{在风场重建中怎样安排这些改进}
表\ref{tab:pinn-optimization-guide}将训练现象与可以优先尝试的方法对应起来。先观察哪一部分学得不好，再选择改进方式，实验过程会更清楚。

\begin{table}[htbp]
\centering\small
\caption{风场 PINN 中常见的训练现象与改进思路。这里给出方法选择的教学对应关系。}
\label{tab:pinn-optimization-guide}
\begin{tabular}{m{0.29\linewidth}m{0.29\linewidth}m{0.30\linewidth}}
\toprule
观察到的现象 & 可以优先检查或尝试 & 主要调整对象\\
\midrule
各项损失量级相差很大 & 无量纲化、残差缩放 & 变量与方程的数值尺度\\
观测拟合较好，方程项下降缓慢 & 梯度监测、自适应加权 & 各项损失参与更新的程度\\
局部高速区或剪切被抹平 & 局部增加配点、傅里叶特征 & 检查位置与细节表达\\
短窗口较好，长窗口变差 & 时间分窗、因果加权 & 单次任务跨度与时间顺序\\
目标方程从随机初始化开始很难训练 & 课程训练、沿用较易任务的参数 & 训练起点与难度安排\\
区域很大，不同位置复杂程度不同 & XPINNs、局部网络与配点配置 & 空间分工与界面连接\\
\bottomrule
\end{tabular}
\end{table}

开始实现时，可以采用如下顺序：
\begin{enumerate}
 \item 先固定一套基础网络，检查坐标、投影和方程实现，统一变量尺度，记录每项损失。
 \item 使用随机配点和 Adam 训练。学习率可以逐步衰减，让训练前期较快调整、后期细化。Wang 等的训练采用 Adam 与指数衰减学习率\cite{wang2023expert}。
 \item 根据结果加入一种有针对性的改进，例如自适应权重或傅里叶特征，并比较相同观测条件下的重建图、留出数据误差和训练耗时。
 \item 当时间跨度或区域规模成为主要困难时，再引入分窗、因果加权或 XPINNs，继续检查窗口交界或子区域界面附近的结果。
\end{enumerate}

还可以沿另外两条方向继续扩展方法。其一是改变方程的表达方式，例如弱形式 PINN 先对方程积分，再形成训练误差，适当的分部积分能够减少直接计算高阶导数的需求\cite{kharazmi2021hpvpinn}。其二是把未知物理参数与风场一起学习，例如估计所选模型中的黏度参数。贝叶斯 PINN 则进一步用概率分布表达参数和预测的不确定程度\cite{yang2021bpinn}。这些扩展可以与前面的训练优化方法结合，用于不同的风场重建任务。

\FloatBarrier

\paragraph{本章小结。}
PINN 的训练同时利用测量值和方程：测量值约束已测位置的预测，方程约束风场在空间与时间中的变化。实现时，先明确网络的输入和输出，再写观测投影、计算方程残差，把它们合成损失。训练完成后，在指定网格和时刻查询网络，就能得到风速分量、风速大小与风向。论文中的双波束案例把这些步骤串在了一起，也展示了从时间窗内的重建到短时外推的使用过程。在此基础上，还可以从变量尺度、损失平衡、时间安排、网络结构和区域分解等方面改善训练。
